\documentclass[11pt, aspectratio=169, xcolor={dvipsnames}, usepdftitle=false]{beamer}
\usepackage{presentation}

\usepackage{amsmath, amsthm, amssymb} \renewcommand{\qedsymbol}{\textsc{qed}}
\usepackage{mathtools}
\usepackage{cmll}
\usepackage{ebproof}
\usepackage{csquotes}
\usepackage{enumitem}

\usepackage{circuitikz}
\usepackage{tikz-cd}
\usetikzlibrary{babel}
% command for labels in commuting diagrams
\newcommand{\zb}[2][1em]{\makebox[#1]{$\displaystyle#2$}}

% Shorthands for (logical) frameworks
\newcommand{\FD}{\textsf{FD}}
\newcommand{\ND}{\textsf{ND}}
\newcommand{\LK}{\textsf{LK}}
\newcommand{\LFD}{\textsf{LFD}}
\newcommand{\LL}{\textsf{LL}}
\newcommand{\UCL}{\textsf{UCL}}
\newcommand{\CP}{\textsf{CP}}
\newcommand{\piFD}{\pi\textsf{FD}}

% Syntax of piFD
% statements for processes in piFD
\newcommand{\seq}{;\,}
\newcommand{\link}[2]{#1 \leftrightarrow #2} \newcommand{\linkbf}[2]{#1 \boldsymbol{\leftrightarrow} #2}
\newcommand{\parcomp}[4]{(\nu #1 #2)(#3 \mid #4)} \newcommand{\parcompbf}[4]{(\boldsymbol{\nu} #1 #2)(#3 \mid #4)}
\newcommand{\lselect}[2]{#1 \lhd inl(#2)}
\newcommand{\rselect}[2]{#1 \lhd inr(#2)}
\newcommand{\choice}[3]{#1 \rhd \{ #2,\; #3 \}}
\newcommand{\send}[3]{#1 ! [#2 \mid #3]}
\newcommand{\receive}[4]{#1 ? (#2, #3).#4}
\newcommand{\close}[1]{#1 ! [\,]}
\newcommand{\wait}[2]{#1 ? ().#2}

\newcommand{\slselect}[1]{\lhd inl(#1)}
\newcommand{\srselect}[1]{\lhd inr(#1)}
\newcommand{\schoice}[2]{\rhd \{ #1,\; #2 \}}
\newcommand{\ssend}[2]{! [#1 \mid #2]} \newcommand{\ssendbf}[2]{\textbf{!} [#1 \mid #2]}
\newcommand{\sreceive}[3]{? (#1, #2).#3}
\newcommand{\sclose}{! [\,]}
\newcommand{\swait}[1]{? ().#1}

\newcommand{\future}[1]{#1}

% Operational semantics
\newcommand{\contr}{\rhd_1}
\newcommand{\red}{\rhd}

\newcommand{\parred}{\twoheadrightarrow}
\newcommand{\simplred}{\unrhd}
\newcommand{\dircong}{\Rrightarrow}

\newcommand{\fv}{\mathsf{fv}}
\newcommand{\dom}{\mathsf{dom}}
\newcommand{\rng}{\mathsf{rng}}
\newcommand{\spine}{\mathsf{spine}}

% Rocq Icon
\newcommand{\rocqicon}{%
  \raisebox{-0.15\baselineskip}{%
    \includegraphics[height=0.85em]{../paper/gfx/rocq-logo.pdf}%
  }%
}

\newcommand{\bool}{\textsf{bool}}
\newcommand{\cobool}{\textsf{bool}^\bot}

\hypersetup{pdfpagemode=UseNone, hidelinks, pdfdisplaydoctitle=true,%
%
% Enter presentation title to populate PDF metadata:
pdftitle={presentation}}

% Enter path to PDF file with figures:
\newcommand{\pdf}{figures.pdf}

% Enter title:
\title{Free Deduction as Sessions}

% Enter authors:
\author{Filip Šimić}

% Enter location and date:
\date{16.09.2026}

\begin{document}
\frame{\titlepage}

% Free Deduction as Sessions
\begin{frame}
  \begin{center}
  \begin{minipage}[t]{0.4\textwidth}
    \centering
    \textbf{Logical Framework}

    \bigskip

    {\LARGE Free Deduction}

    \bigskip

    \textbf{LFD}

    \bigskip

    \begin{flushleft}
    \scalebox{.7}{
    \begin{minipage}[t]{\textwidth}
      \visible<2->{
      \begin{prooftree}[separation=0.7em]
        \hypo{\tau, \tau^\bot \vdash }

        \hypo{A \otimes B, A^\bot \parr B^\bot \vdash}

        \hypo{B \otimes A, B^\bot \parr A^\bot \vdash}
        \hypo{B, B^\bot \vdash}
        \hypo{A, A^\bot \vdash}
        \infer3[$\otimes$]{A, B, B^\bot \parr A^\bot \vdash}

        \infer2[$\parr$]{A \otimes B, B^\bot \parr A^\bot \vdash}

        \infer2[$\multimap$]{(A \otimes B) \prec (B \otimes A) \vdash}
      \end{prooftree}
      }
    \end{minipage}
    }
  \end{flushleft}
  \end{minipage}\hfill
  \begin{minipage}[t]{0.05\textwidth}
    \centering
    \phantom{FOO}

    \bigskip

    \large as

    \bigskip

    \phantom{foo}
  \end{minipage}\hfill
  \begin{minipage}[t]{0.4\textwidth}
    \centering
    \textbf{Process Calculus}

    \bigskip

    {\LARGE Sessions}

    \bigskip

    $\boldsymbol{\pi}$\textbf{FD}

    \bigskip

    \begin{center}
      \scalebox{.7}{
        \begin{minipage}[t]{\textwidth}
          \vspace{-4em}
          \visible<3->{
          \begin{align*}
            &\textbf{swap}(c : (A \otimes B) \prec (B \otimes A)) \coloneq \\
            &\qquad \linkbf{c}{c^\bot} \seq \\
            &\qquad c^\bot \textbf{?} (pair, output) . \{ \\
            &\qquad \qquad pair \boldsymbol{\leftrightarrow} \textbf{?}(a, b) . \{ \\
            &\qquad \qquad \qquad \linkbf{output}{\ssendbf{b^\bot.\linkbf{b}{b^\bot}}{a^\bot.\linkbf{a}{a^\bot}}} \\
            &\qquad \qquad \} \\
            &\qquad \}
          \end{align*}
          }
        \end{minipage}
    }
    \end{center}
  \end{minipage}
  \end{center}
\end{frame}

%%%%%%%%%%%%%%%%%%%%%%%%%%%%%%%%%%%%%%%%%%%%%%%%%%%%%%%%%%%%%%%%%%%%%%%%%%%%%%%%
% Around Introduction and Elimination
%%%%%%%%%%%%%%%%%%%%%%%%%%%%%%%%%%%%%%%%%%%%%%%%%%%%%%%%%%%%%%%%%%%%%%%%%%%%%%%%
\begin{frame}
\heading{Around Introduction and Elimination}
\end{frame}

% Paper slide: This thesis is in a sense a follow up to the Intro/Elim paper
\begin{frame}
  \includegraphics[page=2, trim={0cm 14cm 0cm 3cm}, clip, width=0.9\linewidth]{paper/ostermann-introelim.pdf}
\end{frame}

% Discuss the design matrix for sequent-style calculi
\begin{frame}
  \frametitle{Designing Sequent-Style Calculi}
  \begin{center}
    \begin{minipage}{0.49\textwidth}
      \centering
      \begin{prooftree}
        \hypo{P_1 \hdots P_n}
        \infer1[\textsc{l-$\odot$-intro}]{\Gamma, \odot(T_1, \hdots, T_n) \vdash \Delta}
      \end{prooftree}
      \\ \bigskip \bigskip
      \begin{prooftree}
        \hypo{\Gamma, \odot(T_1, \hdots, T_n) \vdash \Delta}
        \hypo{P_1 \hdots P_n}
        \infer2[\textsc{l-$\odot$-elim}]{C}
      \end{prooftree}
    \end{minipage}\hfill
    \begin{minipage}{0.49\textwidth}
      \centering
      \begin{prooftree}
        \hypo{P_1 \hdots P_n}
        \infer1[\textsc{r-$\odot$-intro}]{\Gamma \vdash \odot(T_1, \hdots, T_n), \Delta}
      \end{prooftree}
      \\ \bigskip \bigskip
      \begin{prooftree}
        \hypo{\Gamma \vdash \odot(T_1, \hdots, T_n), \Delta}
        \hypo{P_1 \hdots P_n}
        \infer2[\textsc{r-$\odot$-elim}]{C}
      \end{prooftree}
    \end{minipage}
  \end{center}
\end{frame}

% Since this work is in the domain of Curry-Howard interpretations
% we can choose to start with Free Deduction or piFD, but we choose the latter

%%%%%%%%%%%%%%%%%%%%%%%%%%%%%%%%%%%%%%%%%%%%%%%%%%%%%%%%%%%%%%%%%%%%%%%%%%%%%%%%
% Programming with Sessions
%%%%%%%%%%%%%%%%%%%%%%%%%%%%%%%%%%%%%%%%%%%%%%%%%%%%%%%%%%%%%%%%%%%%%%%%%%%%%%%%
\begin{frame}
\heading{Programming with Sessions}
\end{frame}

\begin{frame}
  \begin{minipage}[t]{0.45\textwidth}
    \begin{align*}
      &\textbf{not}(client : \bool \prec \bool) \coloneq \\
      \onslide<2,5->{&\qquad \link{client}{server} \seq} \\
      \onslide<3,5->{&\qquad server \textbf{?} (input, output).}\onslide<3,5->{\{} \\
      \onslide<4,5->{&\qquad \qquad \textbf{if} \; input \; \\
      &\qquad \qquad \textbf{then} \; \{ \linkbf{output}{\textbf{false}} \} \\
      &\qquad \qquad \textbf{else} \; \{ \linkbf{output}{\textbf{true}} \}} \\
      \onslide<3,5->{&\qquad \}}
    \end{align*}
  \end{minipage}
  \visible<5->{\begin{minipage}[t]{0.45\textwidth}
    \begin{align*}
      &\textbf{true\_client}(server : \bool \multimap \bool) \coloneq \\
      \onslide<6,8->{&\qquad \linkbf{server}{client} \seq} \\
      \onslide<7->{&\qquad client \textbf{!} [v . \linkbf{v}{\textbf{true}} \mid res. \link{res}{\star}]}
    \end{align*}
  \end{minipage}
  }

  \visible<8->{
  \begin{align*}
    \parcompbf{\;client}{\;server}{\textbf{not}(client)}{\textbf{true\_client}(server)}
  \end{align*}
  }
\end{frame}

\begin{frame}
  \begin{align*}
         &\; \parcompbf{\;client}{\;server}{\textbf{not}(client)}{\textbf{true\_client}(server)} \\
    \red &\; \ssendbf{v . \linkbf{v}{\textbf{true}}}{res . \linkbf{res}{\star}} \leftrightarrow \\
         &\; \textbf{?}(input, output). \\
         &\; \quad \textbf{if} \; input \; \textbf{then} \; \{\linkbf{output}{\textbf{false}}\} \; \textbf{else} \; \{\linkbf{output}{\textbf{true}}\} \\
    \visible<2->{
    \red &\; (\boldsymbol{\nu} \; output \; res) ((\boldsymbol{\nu} \; input \; v) \\
         &\; \quad (\textbf{if} \; input \; \textbf{then} \; \{\linkbf{output}{\textbf{false}}\} \; \textbf{else} \; \{\linkbf{output}{\textbf{true}}\} \\
         &\; \quad \mid \linkbf{v}{\textbf{true}}) \\
         &\; \quad \mid \linkbf{res}{\star}) \\
    }
    \visible<3->{
    \red &\; \textbf{if} \; \textbf{true} \; \textbf{then} \; \{\linkbf{\star}{\textbf{false}}\} \; \textbf{else} \; \{\linkbf{\star}{\textbf{true}}\} \\
    }
    \visible<4->{
    \red &\; \linkbf{\star}{\textbf{false}}
    }
  \end{align*}
\end{frame}

\begin{frame}
\heading{A Tour of Free Deduction}
\end{frame}

\begin{frame}
  \includegraphics[page=3, trim={1cm 13.5cm 1cm 2cm}, clip, width=0.9\linewidth]{paper/parigot-freededuction.pdf}
\end{frame}

\begin{frame}
  \frametitle{The Justifications of the Elimination Rules}
  \begin{center}
    \begin{prooftree}
      \hypo{\Gamma_1 \vdash A, \Delta_1}
      \hypo{\Gamma_2 \vdash B, \Delta_2}
      \infer2[]{\Gamma_1, \Gamma_2 \vdash A \otimes B, \Delta_1, \Delta_2}
    \end{prooftree}
  \end{center}

  \begin{center}
  \begin{minipage}[t]{0.45\textwidth}
    \visible<2->{
    \begin{center}
      \textbf{Inversion Principle}
    \end{center}
    }
    \visible<3->{
    \begin{center}
    \enquote{Whatever follows from the grounds for deriving a
    formula must follow from that formula.}

    \bigskip

      \begin{prooftree}
        \hypo{\Gamma_1 \vdash A \otimes B, \Delta_1}
        \hypo{\Gamma_2, A, B \vdash \Delta_2}
        \infer2[]{\Gamma_1, \Gamma_2 \vdash \Delta_1, \Delta_2}
      \end{prooftree}
    \end{center}
    }
  \end{minipage}
  \begin{minipage}{0.05\textwidth}
    %
  \end{minipage}
  \begin{minipage}[t]{0.45\textwidth}
    \visible<2->{
    \begin{center}
      \textbf{Dual Inversion Principle}
    \end{center}
    }
    \visible<4->{
    \begin{center}
    \enquote{Whatever follows from a formula must follow from
    the sufficient grounds for deriving the formula.}

    \bigskip

      \begin{prooftree}
        \hypo{\Gamma_1, A \otimes B \vdash \Delta_1}
        \hypo{\Gamma_2 \vdash A, \Delta_2}
        \hypo{\Gamma_3 \vdash B, \Delta_3}
        \infer3[]{\Gamma_1, \Gamma_2, \Gamma_3 \vdash \Delta_1, \Delta_2, \Delta_3}
      \end{prooftree}
    \end{center}
    }
  \end{minipage}
  \end{center}
\end{frame}

\begin{frame}
  \frametitle{Cut is Derivable}
  \begin{center}
    \begin{prooftree}
      \hypo{0 \oplus A \vdash 0 \oplus A}
      \hypo{\al{\Gamma \vdash A, \Delta}}
      \infer2[\textsc{l$\oplus$}]{\Gamma \vdash 0 \oplus A, \Delta}

      \hypo{0 \vdash 0}
      \infer1[\textsc{r$0$}]{0 \vdash}
      \hypo{\al{\Pi, A \vdash \Theta}}
      \infer3[\textsc{r$\oplus$}]{\Gamma, \Pi \vdash \Delta, \Theta}
    \end{prooftree}

    \bigskip \bigskip

    What is a \al{cut-free} proof in $\FD$?

    \bigskip

    \visible<2->{
    OR

    \bigskip

    What proofs possess the \al{subformula property}?
    }
  \end{center}
\end{frame}

%%%%%%%%%%%%%%%%%%%%%%%%%%%%%%%%%%%%%%%%%%%%%%%%%%%%%%%%%%%%%%%%%%%%%%%%%%%%%%%%
% Cut Elimination in Free Deduction
%%%%%%%%%%%%%%%%%%%%%%%%%%%%%%%%%%%%%%%%%%%%%%%%%%%%%%%%%%%%%%%%%%%%%%%%%%%%%%%%
\begin{frame}
  \heading{Normalization in Free Deduction}
\end{frame}

% Formula Property
\begin{frame}
  \frametitle{When do Free Deduction Proofs violate the Subformula Property?}
  \textbf{Observation: }
  Every instance of the logical rules removes the principal formula from the sequent entirely.

  \bigskip

  \begin{lemma}[Formula Property]
    Every formula $A$ which occurs in the conclusion $\Gamma \vdash \Delta$
    of a proof in $\FD$ comes from an axiom $A \vdash A$.
  \end{lemma}
\end{frame}

% Logical Redexes
\begin{frame}
  \frametitle{Detours in Proofs of Free Deduction: Logical Redexes}
  To violate the subformula property, both formulas of $A \vdash A$ must be subject of an elimination.

  \bigskip

  \begin{center}
    \begin{prooftree}
      \hypo{}
      \infer1[\textsc{ax}]{A \vdash A}

      \hypo{\mathcal{D}_i}
      \ellipsis{}{\Pi_i \vdash \Theta_i}
      \delims{ \left( }{ \right)_{i} }
      \hypo{\hspace{-1.5em}} % by default, spacing between hypos is 1.5em

      \infer3[\textsc{l$\odot$}]{\Gamma_1 \vdash A, \Delta_1}

      \hypo{\mathcal{D}_k'}
      \ellipsis{}{\Pi'_k \vdash \Theta'_k}
      \delims{ \left( }{ \right)_{k} }
      \hypo{\hspace{-1.5em}}

      \infer3[\textsc{r$\odot$}]{\Gamma_1, \Gamma_2 \vdash \Delta_1, \Delta_2}
    \end{prooftree}
  \end{center}
\end{frame}

% Logical Redex Example: Plus
% \begin{frame}
%   \frametitle{Logical Redex: $\oplus$}
%   \begin{center}
%   \begin{prooftree}
%     \hypo{}
%     \infer1[\textsc{ax}]{A \oplus B \vdash A \oplus B}
%     \hypo{\mathcal{D}_1}
%     \ellipsis{}{\Gamma, A \vdash \Delta}
%     \hypo{\mathcal{D}_2}
%     \ellipsis{}{\Gamma, B \vdash \Delta}
%     \infer3[\textsc{r$\oplus$}]{\Gamma, A \oplus B \vdash \Delta}
%     \hypo{\mathcal{D}_3}
%     \ellipsis{}{\Pi \vdash A, \Theta}
%     \infer2[\textsc{l$\oplus_1$}]{\Gamma, \Pi \vdash \Delta, \Theta}
%   \end{prooftree}
% \end{center}
% \visible<2->{
% may reduce to
% \begin{center}
%   \begin{prooftree}
%     \hypo{\mathcal{D}_1[\mathcal{D}_3 / A]}
%     \ellipsis{}{\Gamma, \Pi \vdash \Delta, \Theta}
%   \end{prooftree}
%   \qquad or \qquad
%   \begin{prooftree}
%     \hypo{\mathcal{D}_3[\mathcal{D}_1 / A]}
%     \ellipsis{}{\Gamma, \Pi \vdash \Delta, \Theta}
%   \end{prooftree}
% \end{center}
% }
% \end{frame}

% Structural Redexes
\begin{frame}
  \frametitle{Detours in Proofs of Free Deduction: Structural Redexes}

  \begin{center}
    \begin{prooftree}
      \hypo{P_1, \hdots, P_{i - 1}}
      \hypo{\mathcal{D}}
      \ellipsis{}{\Gamma_i, A \odot B \vdash \Delta_i}
      \hypo{P_{i + 1}, \hdots, P_n}
      \infer3[$\mathcal{R}$]{\Gamma, A \odot B \vdash \Delta}
      \hypo{P_1', \hdots, P_m'}
      \infer2[\textsc{l$\odot$}]{\Gamma' \vdash \Delta'}
    \end{prooftree}
  \end{center}

  \visible<2->{
  \begin{center}
    $\rotatebox[origin=c]{-90}{$\rightsquigarrow$}$

    \begin{prooftree}
      \hypo{P_1, \hdots, P_{i - 1}}

      \hypo{\mathcal{D}}
      \ellipsis{}{\Gamma_i, A \odot B \vdash \Delta_i}
      \hypo{P_1', \hdots, P_m'}
      \infer2[\textsc{l$\odot$}]{\Gamma'' \vdash \Delta''}

      \hypo{P_{i + 1}, \hdots, P_n}
      \infer3[$\mathcal{R}$]{\Gamma' \vdash \Delta'}
    \end{prooftree}
  \end{center}
  }
\end{frame}

% Cut Elimination and Process Reduction
\begin{frame}
  \frametitle{How do Cut Elimination and Process Reduction relate? - A first Approximation}

  \begin{center}
  \begin{minipage}[t]{0.45\textwidth}
    \begin{center}
      \textbf{Logical Redexes}
    \end{center}
    \begin{center}
    Communication between channel endpoints
    \end{center}
  \end{minipage}
  \begin{minipage}{0.05\textwidth}
    %
  \end{minipage}
  \begin{minipage}[t]{0.45\textwidth}
    \begin{center}
      \textbf{Structural Redexes}
    \end{center}
    \begin{center}
    Fulfilling a promise with a value
    \end{center}
  \end{minipage}
  \end{center}
\end{frame}

%%%%%%%%%%%%%%%%%%%%%%%%%%%%%%%%%%%%%%%%%%%%%%%%%%%%%%%%%%%%%%%%%%%%%%%%%%%%%%%%
% From FD to LFD
%%%%%%%%%%%%%%%%%%%%%%%%%%%%%%%%%%%%%%%%%%%%%%%%%%%%%%%%%%%%%%%%%%%%%%%%%%%%%%%%
\begin{frame}
\heading{From FD to LFD}
\end{frame}

% Linear Negation / Collapsing Chirality
\begin{frame}
  \frametitle{Communication occurs between Dual Types}
  \begin{definition}[Linear Negation]
    \vspace{-1.5em}
    \begin{align*}
      1^\bot &\coloneq \bot
      & \bot^\bot &\coloneq 1 \\
      (A \otimes B)^\bot &\coloneq A^\bot \parr B^\bot
      & (A \parr B)^\bot &\coloneq A^\bot \otimes B^\bot \\
      (A \oplus B)^\bot &\coloneq A^\bot \with B^\bot
      & (A \with B)^\bot &\coloneq A^\bot \oplus B^\bot
    \end{align*}
  \end{definition}

  \bigskip

  \begin{center}
    \begin{prooftree}
      \hypo{}
      \infer1[]{A \vdash A}
    \end{prooftree}
    $\quad \rightsquigarrow \quad$
    \begin{prooftree}
      \hypo{}
      \infer1[]{A, A^\bot \vdash}
    \end{prooftree}
    \visible<2->{
    $\quad \rightsquigarrow \quad$
    \begin{prooftree}
      \hypo{}
      \infer1[]{\al{x} : A, \al{y} : A^\bot \vdash \al{\link{x}{y}}}
    \end{prooftree}
    }
  \end{center}
\end{frame}

% Demonstrate term assignment using tensor
\begin{frame}
  \frametitle{Adding Term Assignments}
  $\qquad$
  \begin{center}
    \begin{prooftree}
      \hypo{\Gamma_1, A \otimes B \vdash \Delta_1}
      \hypo{\Gamma_2 \vdash A, \Delta_2}
      \hypo{\Gamma_3 \vdash B, \Delta_3}
      \infer3[]{\Gamma_1, \Gamma_2, \Gamma_3 \vdash \Delta_1, \Delta_2, \Delta_3}
    \end{prooftree}

    \bigskip\bigskip

    \visible<2->{
    $\rightsquigarrow \qquad$
    \begin{prooftree}
      \hypo{\Gamma_1, A \otimes B \vdash}
      \hypo{\Gamma_2, A^\bot \vdash}
      \hypo{\Gamma_3, B^\bot \vdash}
      \infer3[$\otimes$]{\Gamma_1, \Gamma_2, \Gamma_3 \vdash}
    \end{prooftree}
    $\qquad$
    }

    \bigskip\bigskip

    \visible<3->{
    $\rightsquigarrow \qquad$
    \begin{prooftree}
      \hypo{\Gamma_1, \al{x} : A \otimes B \vdash \al{P}}
      \hypo{\Gamma_2, \al{v} : A^\bot \vdash \al{Q}}
      \hypo{\Gamma_3, \al{z} : B^\bot \vdash \al{R}}
      \infer3[]{\Gamma_1, \Gamma_2, \Gamma_3 \vdash \al{P \seq x \ssend{v.Q}{z.R}}}
    \end{prooftree}
    $\qquad$
    }
  \end{center}
\end{frame}

% Introdouce new syntactic category: Messages
\begin{frame}
  \frametitle{How to Eliminate Structural Redexes in $\pi$FD?}
  \begin{center}
  $\quad\qquad$
  \begin{prooftree}
    \hypo{\Gamma_1, \al{x} : A \otimes B \vdash \al{P}}
    \hypo{\Gamma_2, \al{v} : A^\bot \vdash \al{Q}}
    \hypo{\Gamma_3, \al{z} : B^\bot \vdash \al{R}}
    \infer3[]{\Gamma_1, \Gamma_2, \Gamma_3 \vdash \al{P \seq x \ssend{v.Q}{z.R}}}
  \end{prooftree}
  $\qquad$

  \bigskip\bigskip

  \visible<2->{
  $\rightsquigarrow \qquad$
  \begin{prooftree}
    \hypo{\Gamma_1, \al{x} : A \otimes B \vdash \al{P}}
    \hypo{\Gamma_2, \al{v} : A^\bot \vdash \al{Q}}
    \hypo{\Gamma_3, \al{z} : B^\bot \vdash \al{R}}
    \infer2[\textsc{t$\otimes$}]{\Gamma_2, \Gamma_3 \vdash \al{\ssend{v.Q}{z.R}}}
    \infer2[\textsc{t-seq}]{\Gamma_1, \Gamma_2, \Gamma_3 \vdash \al{P \seq x \ssend{v.Q}{z.R}}}
  \end{prooftree}
  }

  \bigskip\bigskip

  \visible<3->{
  $\contr \qquad P[\al{\ssend{v.Q}{z.R}} \; / \; x]$
  }
  \end{center}
\end{frame}

% Demonstrate missing reduction from introduction
\begin{frame}
  \begin{minipage}[t]{0.45\textwidth}
    \begin{align*}
      &\textbf{true\_client}(server : \bool \multimap \bool) \coloneq \\
      &\qquad \linkbf{server}{client} \seq \\
      &\qquad client \textbf{!} [v . \linkbf{v}{\textbf{true}} \mid res. \link{res}{\star}]
    \end{align*}
  \end{minipage}
  \visible<2->{
  \begin{minipage}[t]{0.04\textwidth}

    \bigskip\bigskip\bigskip

    $\rightsquigarrow$
  \end{minipage}
  \begin{minipage}[t]{0.45\textwidth}
    \begin{align*}
      &\textbf{true\_client}(server : \bool \multimap \bool) \coloneq \\
      &\qquad server \boldsymbol{\leftrightarrow} \, \textbf{!} [v . \linkbf{v}{\textbf{true}} \mid res. \link{res}{\star}]
    \end{align*}
  \end{minipage}
  }
\end{frame}

\begin{frame}
  \frametitle{Syntax}

  \begin{definition}[Syntax of $\piFD$]
  \phantom{foo}
  % \vspace{-4em}
  \begin{minipage}[t]{0.45\textwidth}
    \vspace{-2em}
    \begin{flalign*}
      P, Q \; &\Coloneq &\!\!\!\!\!&                      & \textbf{Process} \\
              &         &\!\!\!\!\!& \link{M}{M}          & \textit{link} \\
              &\;\;\mid &\!\!\!\!\!& P \seq xS            & \textit{sequencing} \\
              &\;\;\mid &\!\!\!\!\!& \Lambda              & \textit{nil process} \\
              \visible<2->{&\;\;\mid &\!\!\!\!\!&  \al{\parcomp{x}{y}{P}{Q}} & \al{\textit{cut}}} \\
      M \;    &\Coloneq &\!\!\!\!\!&                      & \textbf{Message} \\
              &         &\!\!\!\!\!& \future{x}           & \textit{future} \\
              &\;\;\mid &\!\!\!\!\!& S                    & \textit{statement}
    \end{flalign*}
  \end{minipage}
  \hfill
  \begin{minipage}[t]{0.45\textwidth}
    \vspace{-2em}
    \begin{flalign*}
      S \; &\Coloneq &\!\!\!\!\!&                    & \textbf{Statement} \\
           &         &\!\!\!\!\!& \slselect{x.P}     & \textit{selection} \\
           &\;\;\mid &\!\!\!\!\!& \srselect{x.P}     & \textit{selection} \\
           &\;\;\mid &\!\!\!\!\!& \schoice{x.P}{x.P} & \textit{choice} \\
           &\;\;\mid &\!\!\!\!\!& \ssend{x.P}{x.P}   & \textit{output} \\
           &\;\;\mid &\!\!\!\!\!& \sreceive{x}{x}{P} & \textit{input} \\
           &\;\;\mid &\!\!\!\!\!& \sclose            & \textit{close} \\
           &\;\;\mid &\!\!\!\!\!& \swait{P}          & \textit{wait}
    \end{flalign*}
  \end{minipage}
\end{definition}
\end{frame}

\begin{frame}
  \frametitle{Cut enables Communication between Concurrent Processes}
  \begin{center}
    \begin{prooftree}
      \hypo{\Gamma, A \vdash}
      \hypo{\Delta, A^\bot \vdash}
      \infer2[]{\Gamma, \Delta \vdash}
    \end{prooftree}
    $\quad\rightsquigarrow\quad$
    \begin{prooftree}
      \hypo{\Gamma, \al{x} : A \vdash \al{P}}
      \hypo{\Delta, \al{y} : A^\bot \vdash \al{Q}}
      \infer2[]{\Gamma, \Delta \vdash \al{\parcomp{x}{y}{P}{Q}}}
    \end{prooftree}
  \end{center}

  % \visible<2->{
  % \begin{example}
  %   With $A \prec B \coloneq A \otimes B^\bot$ and $A \multimap B \coloneq A^\bot \parr B$,
  %   we have

  %   \begin{center}
  %     \begin{prooftree}
  %       \hypo{client : \bool \prec \bool \vdash \textbf{not}(client)}
  %       \hypo{server : \bool \multimap \bool, \star : \bool^\bot \vdash \textbf{true\_client}(server)}
  %       \infer2[]{\star : \bool^\bot \vdash \parcompbf{\;client}{\;server}{\textbf{not}(client)}{\textbf{true\_client}(server)}}
  %     \end{prooftree}
  %   \end{center}
  % \end{example}
  % }
\end{frame}

\begin{frame}
  \frametitle{Cut Reduction as Message Passing between Processes}
  \vspace{-2em}
  \begin{align*}
    \parcomp{x}{y}{\link{x}{M}}{P}
    &\contr P[M/y]
    \tag{\textsc{r-axcut}}
  \end{align*}

  \visible<3->{
  \begin{center}
    \begin{prooftree}
      \hypo{P \contr P'}
      \infer1[\textsc{r-cong-cut}]{\parcomp{x}{y}{P}{Q} \contr \parcomp{x}{y}{P'}{Q}}
    \end{prooftree}
    \quad
    \begin{prooftree}
      \hypo{P \equiv P'}
      \hypo{P' \contr Q'}
      \hypo{Q' \equiv Q}
      \infer3[\textsc{r-struct}]{P \contr Q}
    \end{prooftree}
  \end{center}
  }

  \visible<2->{
  \begin{example}
    \vspace{-2em}
    \begin{align*}
           &\; \parcompbf{\;client}{\;server}{\textbf{not}(client)}{\textbf{true\_client}(server)} \\
         = &\; \parcompbf{\;client}{\;server}{\textbf{not}(client)}{server \boldsymbol{\leftrightarrow} \, \textbf{!} [v . \linkbf{v}{\textbf{true}} \mid res. \link{res}{\star}]} \\
      \red &\; \ssendbf{v . \linkbf{v}{\textbf{true}}}{res . \linkbf{res}{\star}} \leftrightarrow \\
           &\; \textbf{?}(input, output). \\
           &\; \quad \textbf{if} \; input \; \textbf{then} \; \{\linkbf{output}{\textbf{false}}\} \; \textbf{else} \; \{\linkbf{output}{\textbf{true}}\}
    \end{align*}
  \end{example}
  }
\end{frame}

\begin{frame}
  \frametitle{Logical Cut Reduction for Input and Output}
  \begin{equation*}
      \begin{aligned}
        &\begin{prooftree}[separation=0.5em]
            \hypo{A \otimes B, A^\bot \parr B^\bot \vdash}
            \hypo{\mathcal{D}_A}
            \ellipsis{}{\Gamma, A^\bot \vdash}
            \hypo{\mathcal{D}_B}
            \ellipsis{}{\Delta, B^\bot \vdash}
            \infer3[$\otimes$]{\Gamma, \Delta, A^\bot \parr B^\bot \vdash}
            \hypo{\mathcal{D}}
            \ellipsis{}{\Theta, A, B \vdash}
            \infer2[$\parr$]{\Gamma, \Delta, \Theta \vdash}
          \end{prooftree}
        % \onslide<1>{
        % \quad\contr\quad
        % \begin{prooftree}
        %   \hypo{\mathcal{D}[\mathcal{D}_A / A, \mathcal{D}_B / B]}
        %   \ellipsis{}{\Gamma, \Delta, \Theta \vdash}
        % \end{prooftree}
        % }
        \\ \\
        % \visible<2->{
        &\qquad\qquad\contr\quad
        \begin{prooftree}[separation=0.5em]
          \hypo{\Gamma, A^\bot \vdash}
          \hypo{\Delta, B^\bot \vdash}
          \hypo{\Theta, A, B \vdash}
          \infer2[\textsc{cut}]{\Delta, \Theta, A \vdash}
          \infer2[\textsc{cut}]{\Gamma, \Delta, \Theta \vdash}
        \end{prooftree}
        % }
      \end{aligned}
  \end{equation*}
  \vspace{-1em}
  % \visible<3->{
  \begin{align}
    % tensor/par
    \link{\ssend{v.P}{z.Q}}{\sreceive{a}{b}{R}}
    &\quad\contr\quad \parcomp{b}{z}{\parcomp{a}{v}{R}{P}}{Q}
    \tag{\textsc{r-$\otimes\parr$}}
  \end{align}
  % }
\end{frame}

\begin{frame}
  \frametitle{Logical Cut Reduction for Input and Output}
  \vspace{-2em}
  \begin{align}
    % tensor/par
    \link{\ssend{v.P}{z.Q}}{\sreceive{a}{b}{R}}
    &\quad\contr\quad \parcomp{b}{z}{\parcomp{a}{v}{R}{P}}{Q}
    \tag{\textsc{r-$\otimes\parr$}}
  \end{align}

  \begin{example}
    \vspace{-2em}
    \begin{align*}
           &\; \ssendbf{v . \linkbf{v}{\textbf{true}}}{res . \linkbf{res}{\star}} \leftrightarrow \\
           &\; \textbf{?}(input, output). \\
           &\; \quad \textbf{if} \; input \; \textbf{then} \; \{\linkbf{output}{\textbf{false}}\} \; \textbf{else} \; \{\linkbf{output}{\textbf{true}}\} \\
      \red &\; (\boldsymbol{\nu} \; output \; res) ((\boldsymbol{\nu} \; input \; v) \\
           &\; \quad (\textbf{if} \; input \; \textbf{then} \; \{\linkbf{output}{\textbf{false}}\} \; \textbf{else} \; \{\linkbf{output}{\textbf{true}}\} \\
           &\; \quad \mid \linkbf{v}{\textbf{true}}) \\
           &\; \quad \mid \linkbf{res}{\star}) \\
    \end{align*}
    \vspace{-4em}
  \end{example}
\end{frame}

\begin{frame}
\heading{Metatheory}
\end{frame}

\begin{frame}
  \frametitle{Type Safety}
  \begin{theorem}[\rocqicon~Preservation]
    If $\Gamma \vdash P$ and $P \contr Q$, then $\Gamma \vdash Q$.
  \end{theorem}

  \begin{theorem}[\rocqicon~Progress]
    If $\Gamma \vdash P$, then $P \contr P'$ for some process $P'$ or
    $P$ is final.
  \end{theorem}

  \begin{corollary}[\rocqicon~Progress for closed processes]
    If $\vdash P$, then $P \contr P'$ for some process $P'$ or
    $P = \Lambda$.
  \end{corollary}

  % \begin{theorem}[\rocqicon~Final processes are irreducible]
  %   If $\Gamma \vdash P$, then $P$ is final if and only if $P$ is irreducible.
  % \end{theorem}
\end{frame}

\begin{frame}[fragile]
  \frametitle{The Church-Rosser Property}
  \begin{theorem}[\rocqicon~Church-Rosser]
    Let $\Gamma \vdash P$ be a well-typed process.
    If $P \red Q$ and $P \red N$, then $Q \red R$ and $N \red R$ for some $R$.
    Pictorially,
    \begin{center}
      \begin{tikzcd}[
          row sep={1.5cm,between origins},
          column sep={1.5cm,between origins}
        ]
            & \zb{P} \arrow[dl, "\red"'] \arrow[dr, "\red"] & \\
        \zb{Q} \arrow[dr, dashrightarrow, "\red"'] &   & \zb{N} \arrow[dl, dashrightarrow, "\red"] \\
            & \zb{R} &
      \end{tikzcd}
    \end{center}
  \end{theorem}
\end{frame}

\begin{frame}
  \frametitle{Conclusion}

  \begin{itemize}
    \item \textbf{Elimination-Based Foundation for Session Calculi}
      \begin{itemize}
        \item Shifted from traditional introduction-rule foundations to an elimination-only framework
      \end{itemize}

    \bigskip

    \item \textbf{Novel Curry-Howard Interpretation for Session-Based Communication}
      \begin{itemize}
        \item Link processes represent the state of a channel.
        \item Communication between concurrent processes (\textsc{cut}) and channel endpoints (links) is disentangled.
      \end{itemize}

    \bigskip

    \item \textbf{Mechanized Metatheory}
      \begin{itemize}
        \item Fully formalized in \rocqicon~Rocq (Preservation, Progress, and Church-Rosser).
      \end{itemize}
  \end{itemize}
\end{frame}

\lastslide

%%%%%%%%%%%%%%%%%%%%%%%%%%%%%%%%%%%%%%%%%%%%%%%%%%%%%%%%%%%%%%%%%%%%%%%%%%%%%%%%
% Backup slides
%%%%%%%%%%%%%%%%%%%%%%%%%%%%%%%%%%%%%%%%%%%%%%%%%%%%%%%%%%%%%%%%%%%%%%%%%%%%%%%%

% Structural Congruence
\begin{frame}
  \frametitle{Structural Congruence: Factoring out Syntactic Differences}
  \begin{definition}[Structural congruence]
    Structural congruence $\equiv$ is defined as the least congruence relation
    that is closed under rules
    \begin{align*}
      \link{x}{y} &\equiv \link{y}{x} \tag{\textsc{c-link}} \\
      \parcomp{x}{y}{P}{Q} &\equiv \parcomp{y}{x}{Q}{P} \tag{\textsc{c-comm}} \\
      \parcomp{v}{z}{\parcomp{x}{y}{P}{Q}}{R}
        &\equiv \parcomp{x}{y}{P}{\parcomp{v}{z}{Q}{R}} \tag{\textsc{c-assoc}}
    \end{align*}
    where for \textsc{c-assoc} we presuppose that $v \notin \fv(P)$ and $y \notin \fv(R)$.
  \end{definition}
\end{frame}

% Typing
\begin{frame}
  \begin{center}
    \begin{flushright}
      Processes \framebox{$\Gamma \vdash P$}
    \end{flushright}

    \begin{prooftree}
      \hypo{\Gamma \vdash M_1 : A}
      \hypo{\Delta \vdash M_2 : A^\bot}
      \infer2[\textsc{t-ax}]{\Gamma \circ \Delta \vdash \link{M_1}{M_2}}
    \end{prooftree}

    \bigskip

    \begin{prooftree}
      \hypo{\Gamma, x : A \vdash P}
      \hypo{\Delta, y : A^\bot \vdash Q}
      \infer2[\textsc{t-cut}]{\Gamma \circ \Delta \vdash \parcomp{x}{y}{P}{Q}}
    \end{prooftree}
    \hfill
    \begin{prooftree}
      \hypo{\Gamma, x : A \vdash P}
      \hypo{\Delta \vdash S : A}
      \infer2[\textsc{t-seq}]{\Gamma \circ \Delta \vdash P \seq xS}
    \end{prooftree}

    \bigskip

    \begin{prooftree}
      \hypo{}
      \infer1[\textsc{t-nil}]{\vdash \Lambda}
    \end{prooftree}
  \end{center}
\end{frame}

\begin{frame}
  \begin{center}
    \begin{flushright}
      Messages \framebox{$\Gamma \vdash M : A$} and
      statements \framebox{$\Gamma \vdash S : A$}
    \end{flushright}
    \vspace{1em}

    \begin{prooftree}[center=false]
      \hypo{}
      \infer1[\textsc{t-id}]{x : A \vdash x : A}
    \end{prooftree}

    \bigskip

    \begin{prooftree}
      \hypo{\Gamma, z : A^\bot \vdash Q}
      \infer1[\textsc{t-$\oplus_1$}]{\Gamma \vdash \slselect{z.Q} : A \oplus B}
    \end{prooftree}
    \qquad
    \begin{prooftree}
      \hypo{\Gamma, z : B^\bot \vdash Q}
      \infer1[\textsc{t-$\oplus_2$}]{\Gamma \vdash \srselect{z.Q} : A \oplus B}
    \end{prooftree}

    \bigskip

    \begin{prooftree}
      \hypo{\Gamma, z_1 : A^\bot \vdash Q_1}
      \hypo{\Gamma, z_2 : B^\bot \vdash Q_2}
      \infer2[\textsc{t-$\with$}]{\Gamma \vdash \schoice{z_1.Q_1}{z_2.Q_2} : A \with B}
    \end{prooftree}

    \bigskip

    \begin{prooftree}
      \hypo{\Gamma, v : A^\bot \vdash Q}
      \hypo{\Delta, z : B^\bot \vdash R}
      \infer2[\textsc{t-$\otimes$}]{\Gamma \circ \Delta \vdash \ssend{v.Q}{z.R} : A \otimes B}
    \end{prooftree}
    \qquad
    \begin{prooftree}
      \hypo{\Gamma, v : A^\bot, z : B^\bot \vdash Q}
      \infer1[\textsc{t-$\parr$}]{\Gamma \vdash \sreceive{v}{z}{Q} : A \parr B}
    \end{prooftree}

    \bigskip

    \begin{prooftree}[center=false]
      \hypo{}
      \infer1[\textsc{t-$1$}]{\vdash \sclose{} : 1}
    \end{prooftree}
    \qquad
    \begin{prooftree}[center=false]
      \hypo{\Gamma \vdash Q}
      \infer1[\textsc{t-$\bot$}]{\Gamma \vdash \swait{Q} : \bot}
    \end{prooftree}
  \end{center}
\end{frame}

% Operational Semantics
\begin{frame}
  \frametitle{Operational Semantics}
  \begin{definition}[Contraction $\contr$]
  \begin{align}
    % tensor/par
    \link{\ssend{v.P}{z.Q}}{\sreceive{a}{b}{R}}
    &\contr \parcomp{b}{z}{\parcomp{a}{v}{R}{P}}{Q}
    \tag{\textsc{r-$\otimes\parr$}}
    \\
    % plus/with
    \link{\slselect{z.P}}{\schoice{v.Q}{w.R}}
    &\contr \parcomp{z}{v}{P}{Q}
    \tag{\textsc{r-$\oplus\with_1$}}
    \\
    \link{\srselect{z.P}}{\schoice{v.Q}{w.R}}
    &\contr \parcomp{z}{w}{P}{R}
    \tag{\textsc{r-$\oplus\with_2$}}
    \\
    % 1/bot
    \link{\sclose}{\swait{P}}
    &\contr P
    \tag{\textsc{r-$1\bot$}}
    \\
    % AxCut
    \parcomp{x}{y}{\link{x}{M}}{P}
    &\contr P[M/y]
    \tag{\textsc{r-axcut}}
    \\
    % Seq
    P \seq xS
    &\contr P[S/x]
    \tag{\textsc{r-seq}}
  \end{align}
\end{definition}
\end{frame}

\begin{frame}
  \frametitle{Operational Semantics (continued)}
  \begin{definition}[Contraction $\contr$ (continued)]
  \begin{center}
    % congruence: cut
    \hfill
    \begin{prooftree}
      \hypo{P \contr P'}
      \infer1[]{\parcomp{x}{y}{P}{Q} \contr \parcomp{x}{y}{P'}{Q}}
    \end{prooftree}
    \hfill (\textsc{r-cong-cut})

    \bigskip

    % congruence: seq
    \hfill
    \begin{prooftree}
      \hypo{P \contr P'}
      \infer1[]{P \seq s \contr P' \seq s}
    \end{prooftree}
    \hfill
    (\textsc{r-cong-seq})

    \bigskip

    % struct
    \hfill
    \begin{prooftree}
      \hypo{P \equiv P'}
      \hypo{P' \contr Q'}
      \hypo{Q' \equiv Q}
      \infer3[]{P \contr Q}
    \end{prooftree}
    \hfill
    (\textsc{r-struct})
  \end{center}
\end{definition}
\end{frame}

% Ephemeral Booleans
\begin{frame}
  \frametitle{A Church-Encoding for Linear Booleans}
  \vspace{-2em}
  \begin{align*}
    \bool &\coloneq 1 \oplus 1 \\
    \cobool &\coloneq \bot \with \bot \\
    \\[-1em]
    true(x : \cobool) &\coloneq \link{x}{y} \seq y \slselect{z.\link{z}{\sclose}} \\
    false(x : \cobool) &\coloneq \link{x}{y} \seq y \srselect{z.\link{z}{\sclose}} \\
    \\[-1em]
    true &\coloneq \slselect{z.\link{z}{\sclose}} \\
    false &\coloneq \srselect{z.\link{z}{\sclose}} \\
    \\[-1em]
    if \; x \; then \; P \; else \; Q
    &\coloneq \link{x}{\schoice{x.\link{x}{\swait{P}}}{x.\link{x}{\swait{Q}}}}
  \end{align*}
\end{frame}

\begin{frame}
  \frametitle{A Church-Encoding for Linear Booleans (continued)}
  \vspace{-2em}
  \begin{align*}
          &\; if \; x \; then \; P \; else \; Q \seq x \leftarrow true \\
         =&\; \link{x}{\schoice{x.\link{x}{\swait{P}}}{x.\link{x}{\swait{Q}}}} \seq x \leftarrow true \\
    \contr&\; \link{true}{\schoice{x.\link{x}{\swait{P}}}{x.\link{x}{\swait{Q}}}} \\
         =&\; \link{\slselect{z.\link{z}{\sclose}}}{\schoice{x.\link{x}{\swait{P}}}{x.\link{x}{\swait{Q}}}} \\
    \contr&\; \parcomp{z}{x}{\link{z}{\sclose}}{\link{x}{\swait{P}}} \\
    \contr&\; \link{\sclose}{\swait{P}} \\
    \contr&\; P
  \end{align*}
\end{frame}

% Normalforms
\begin{frame}
  \begin{definition}[Weak normal form]
  A derivation $\mathcal{D}$ is in weak normal form, if $\mathcal{D} \equiv \mathcal{N}$,
  where $\mathcal{N}$ is of the form
  \begin{center}
    \begin{prooftree}
      \hypo{\mathcal{D}_1}
      \hypo{\mathcal{D}_2}
      \hypo{\hdots}
      \hypo{\mathcal{D}_{n-1}}
      \ellipsis{}{\Gamma_{n-1}, A_{n-1}^\bot, C_{n-1} \vdash}
      \hypo{\mathcal{D}_n}
      \ellipsis{}{\Gamma_n, A_n^\bot, C_{n - 1}^\bot \vdash}
      \infer2[\textsc{cut}]{\Gamma_{n-1}', A_{n-1}^\bot, A_n^\bot, C_{n-2}^\bot \vdash}
      \infer[no rule]1{\vdots}
      \infer3[\textsc{cut}]{\Gamma_2', A_2^\bot, \hdots, A_n^\bot, C_1^\bot \vdash}
      \infer2[\textsc{cut}]{\Gamma_1', A_1^\bot, A_2^\bot, \hdots, A_n^\bot \vdash}
    \end{prooftree}
  \end{center}
  such that every $\mathcal{D}_i$ has the shape
  \vspace{-1em}
  \begin{center}
    \begin{prooftree}
      \hypo{A_i, A_i^\bot \vdash}
      \hypo{(P_i)^j}
      \infer2[$\odot_i$]{\Gamma_i, A_i^\bot, C_i \vdash}
    \end{prooftree}
  \end{center}
  % \vspace{-1em}
  $\odot_i$ is the principal connective of the formula $A_i$ and $C_i$ is the $i$-th
  \textsc{cut}'s principal formula.
\end{definition}
\end{frame}

\begin{frame}[fragile]
  \frametitle{Which processes do we consider as terminated and non-deadlocked?}
  \begin{definition}[Link list]
    The set of \emph{link lists} and their associated \emph{spines} (sequences of names) are
    inductively defined by the following clauses.
    \begin{enumerate}[label=(\roman*)]
      \item If $P_1 \equiv \link{x_1}{M_1}$ and $P_2 \equiv \link{x_2}{M_2}$ such that
            $M_1$ and $M_2$ are not futures, then $\parcomp{x}{y}{P_1}{P_2}$
            is a link list with spine $\spine(P) = [x_1, x_2]$.
      \item If $L \equiv \link{x}{M}$ such that $M$ is not a future and
            $Q$ is a link list with spine $\spine(Q) = \vec{x}$, then 
            $P = \parcomp{a}{b}{L}{Q}$ is a link list with spine $\spine(P) = x :: \vec{x}$.
    \end{enumerate}
  \end{definition}
\end{frame}

\begin{frame}[fragile]
  \frametitle{Which processes do we consider as terminated and non-deadlocked?}
  \begin{definition}[Final process]
    A process $P$ is \emph{final} if it satisfies one of the following conditions:
    \begin{enumerate}[label=(\roman*)]
      \item $P = \Lambda$.
      \item $P$ is of the form $\link{x}{M}$ or $\link{M}{x}$.
      \item There exists a link list $Q$ such that $P \equiv Q$ and
            every channel $a \in \spine(Q)$ is free in $Q$.
    \end{enumerate}
  \end{definition}
\end{frame}

\end{document}